% D.85. Root, FGH, and PBC full independent proof/source/certificate audits PASS.
\subsection{Every globally primitive rank-two Kummer pair}

\begin{statement}{D.85}{Proved}{The whole compact kernel for every globally independent residual pair}
Let \(s=(a,b)\) be arbitrary completed Kummer classes over
\(\mathbb Q\), with arbitrary infinite ramification, and assume
\[
 \operatorname{span}_{\mathbb F_5}
       \{(a_q,b_q)\bmod5:q\text{ prime}\}=\mathbb F_5^2.
                                                               \tag{D.85.1}
\]
The prime five is included in this hypothesis. Write the actual
local classes at five as
\[
 a=A\kappa(5)+\alpha\kappa(6),\qquad
 b=B\kappa(5)+\beta\kappa(6),
 \qquad \mathbf r=(A,B),\quad\mathbf z=(\alpha,\beta).
                                                               \tag{D.85.2}
\]
Put
\[
 \begin{gathered}
 W=\operatorname{span}_{\mathbb F_5}\{\overline{\mathbf r},
                                      \overline{\mathbf z}\},\\
 \epsilon(W)=
 \begin{cases}
 1,&W=0\text{ or }W\text{ is a line other than }
 \langle(1,0)\rangle,\langle(0,1)\rangle,\langle(1,1)\rangle,\\
 0,&\text{otherwise},
 \end{cases}\\
 n=\min\{v_5(A),v_5(\alpha)\}
  +\min\{v_5(B),v_5(\beta)\}
  +\min\{v_5(A-B),v_5(\alpha-\beta)\},\qquad v_5(0)=\infty.
 \end{gathered}                                                \tag{D.85.3}
\]
For D.63's actual Alexander invariants and D.77's compact
convention \(Z(n)\), the complete answer is
\[
 \mathcal S_s\simeq(\mathbb Z/5)^{\epsilon(W)},\qquad
 \mathcal T_s=0,\qquad
 \mathcal K_s\simeq(\mathbb Z/5)^{\epsilon(W)}\times Z(n).
                                                               \tag{D.85.4}
\]
Thus D.82's classification holds under global residual
independence, without requiring two independent directions
away from five. The splitting in (D.85.4) is continuous and
noncanonical.

The new finite-index case includes the actual point
\[
 t=-799,\qquad
 (a,b)=(\kappa(17)+\kappa(47),\,5\kappa(2)+2\kappa(5)),
 \qquad \mathcal K_s\simeq\mathbb Z_5.
                                                               \tag{D.85.5}
\]
Its metabelian fibre exists and is a torsor under
\(\prod_{\mathbb N}\mathbb Z_5\), by D.74. No existence is
asserted for every pair in (D.85.1). Globally dependent
residual pairs, arbitrary infinitely ramified first-cohomology
ranks, and the first higher-derived obstruction are not
computed by this statement.
\end{statement}

\begin{proof}
Throughout, \(\Lambda=\mathbb Z_5[[t_1-1,t_2-1]]\), and
\[
 f_{\mathbf v}=t_1^{v_1}t_2^{v_2}-1,
 \qquad P=(t_1-1)(t_2-1)(t_1t_2-1).
\]
Use the actual untwisted image \(I_0\), its descended Kummer
cocycle \(k_5\), and \(N_1=\ker\chi\cap\ker k_5\) from
D.82. In particular \(I_0/N_1=\mathbb Z_5(1)\rtimes\Gamma\),
where \(\Gamma=\mathbb Z_5^\times\), is the actual quotient.
The compact action has augmentation \(\chi^2\).
There is an actual element of \(N_1\) whose beta multiplier
minus one is \(P\) times a unit. Every element of \(\ker\chi\)
has beta multiplier congruent to one modulo \(P\).

\paragraph{The only remaining shape of the away-five lattice.}
Let
\[
 L=\overline{\sum_{q\ne5}\mathbb Z_5(a_q,b_q)},
 \qquad
 \mathbf z=\sum_{q\ne5}\lambda_q(a_q,b_q),
 \qquad \lambda_q=\frac{\log(q^4)}{4\log6}.
                                                               \tag{D.85.6}
\]
These are D.83's closed lattice and convergent local unit
vector. If \(L\bmod5\) has rank two, D.82 applies. It
cannot have rank zero under (D.85.1), because the only
remaining vector would be \(\mathbf r\). Suppose therefore
that it has rank one. Then \(\overline{\mathbf r}\) is
independent of that line. Choose \(\mathbf u\in L\) whose
reduction generates the line. The vectors
\(\mathbf u,\mathbf r\) form a basis of \(\mathbb Z_5^2\).

Projection onto the \(\mathbf r\)-coordinate in this basis
sends \(L\) to an ideal contained in \(5\mathbb Z_5\).
If it is zero, then \(L=\mathbb Z_5\mathbf u\), and D.84
applies. Otherwise it is \(5^h\mathbb Z_5\), \(h\ge1\).
Subtracting the \(\mathbf u\)-coordinate of an element
with primitive projected image, and multiplying by a unit,
gives the exact decomposition
\[
 L=\mathbb Z_5\mathbf u\oplus5^h\mathbb Z_5\mathbf r,
 \qquad h\ge1.
                                                               \tag{D.85.7}
\]
Thus it remains to treat this finite-index, nonsaturated
away-five lattice.

Write each prime vector uniquely as
\((a_q,b_q)=c_q\mathbf u+5^h d_q\mathbf r\).
The sequences \((c_q),(d_q)\) are \(c_0\): division by
the fixed scalar \(5^h\) preserves convergence to zero.
Consequently the actual Kummer sums
\[
 c=\sum_{q\ne5}c_q\kappa(q),\quad
 d_0=\sum_{q\ne5}d_q\kappa(q),\quad
 (a,b)=\mathbf u c+\mathbf r(\kappa(5)+5^h d_0)
                                                               \tag{D.85.8}
\]
converge uniformly. Put
\[
 \begin{gathered}
 \lambda=\sum_{q\ne5}\lambda_qc_q,\qquad
 \nu=\sum_{q\ne5}\lambda_qd_q,\qquad
 \mathbf z=\lambda\mathbf u+5^h\nu\mathbf r,\\
 e=v_5(\lambda),\qquad
 d=\max\{v_5(r_1),v_5(r_2),v_5(r_1-r_2)\}.
 \end{gathered}
                                                               \tag{D.85.9}
\]
Here \(d_0\) denotes the Kummer class, while \(d\) is
the boundary depth. The latter is zero unless the primitive
vector \(\mathbf r\) has exactly one boundary form divisible
by five; that form then has depth \(d\), possibly infinite.

\paragraph{The actual images and the natural restriction map.}
D.83's integral gauge is a uniformly convergent sum in
\(L\). Its image computation therefore applies verbatim,
including infinite ramification, and gives
\[
 \begin{split}
 I_s&=\{(\mathbf l+\mathbf r k_5(g),g):
                               \mathbf l\in L,\ g\in I_0\},\\
 I_5&=\{(j\mathbf z+\mathbf r k_5(g),g):
                               j\in\mathbb Z_5,\ g\in I_0\}.
 \end{split}                                                   \tag{D.85.10}
\]
Both retain the same actual graph complement. They are
the faithful compact operator images, not abstract groups
with an assigned action.

Let \(D=A_s^D\) be the discrete Cartier dual and define
\[
 I_L=(f_{\mathbf u},f_{5^h\mathbf r}),\qquad
 B=\Lambda/I_L,\qquad R=\Lambda/(f_{\mathbf z}),\qquad
 J=I_LR.
                                                               \tag{D.85.11}
\]
The defining relation for \(R\) belongs to \(I_L\).
The actual \(I_0\)-action on both coefficient quotients is
\[
             \rho_{\mathbf r}(g)=t^{\mathbf r k_5(g)}\rho_0(g).
                                                               \tag{D.85.12}
\]
Translation coinduction gives natural identifications
\[
 H^1(I_s,D)=H^1(I_0,B^D),\qquad
 H^1(I_5,D)=H^1(I_0,R^D).
                                                               \tag{D.85.13}
\]
For the first identity, one may choose the finite coset
transversal of \(L\) in \(\mathbb Z_5^2\); the continuous
function model for \(D\) is then coinduced for \(L\).
For the second, the line transversal, including nonprimitive
or zero \(\mathbf z\), is the one constructed in D.81.
In both cases the translation group's higher cohomology
vanishes, so the Hochschild--Serre edge map gives (D.85.13).
These function identifications show that restriction is
exactly induced by \(B^D\hookrightarrow R^D\); no unrelated
choice of the two coinduction isomorphisms is made.

The exact sequence dual to \(0\to J\to R\to B\to0\)
therefore identifies the whole restriction kernel with
\[
 \mathcal R_s=(J^D)^{I_0}/\operatorname{im}((R^D)^{I_0}).
                                                               \tag{D.85.14}
\]
Finite affine Chebotarev descent, at each finite coefficient
level as in D.78 and D.83, injects
\(\mathcal S_s^*\) into (D.85.14).
Every global boundary depth is zero by (D.85.1), so
\(D^{I_s}=(B^D)^{I_0}=0\) by D.77. In particular the
map in the denominator of (D.85.14) is injective.

\paragraph{An exact two-generator module calculation.}
The \(\mathbb Z_5\)-basis \(\mathbf u,\mathbf r\) gives
genuine completed group-ring coordinates
\[
 x=f_{\mathbf u},\qquad y=f_{\mathbf r},\qquad
 \Lambda=\mathbb Z_5[[x,y]],\qquad
 g_h=(1+y)^{5^h}-1,
 \quad I_L=(x,g_h).
                                                               \tag{D.85.15}
\]
Let \(V=I_L/yI_L\). There is an exact topological
\(\mathbb Z_5[[x]]\)-module identification
\[
 V=x\mathbb Z_5[[x]]\oplus\mathbb Z_5\,g_h,
 \qquad xg_h=0.
                                                               \tag{D.85.16}
\]
Indeed \(g_h=yS_h(y)\), where \(S_h(0)=5^h\ne0\).
In \(V\), both \(yg_h\) and \(xg_h\) vanish; the latter
belongs to \(yx\Lambda\). Every element thus has the
form in (D.85.16). To prove uniqueness, suppose
\(xf(x)+a g_h\in y(x,g_h)\). Evaluation at \(y=0\)
gives \(f=0\). Cancel \(y\) and set \(x=0\); then
\(aS_h(y)=yS_h(y)b(y)\). Cancellation in the power-series
domain gives \(a=yb(y)\), hence \(a=0\). The map and its
inverse are continuous, since they are continuous maps
between compact Hausdorff modules.

The exact local relation has the form
\[
 f_{\mathbf z}=(1+x)^\lambda(1+g_h)^\nu-1.
                                                               \tag{D.85.17}
\]
Modulo \(yI_L\), the products \(g_h^2\) and \(xg_h\)
vanish, so its image is
\[
                 f_\lambda(x)+\nu g_h,
 \qquad f_\lambda(x)=(1+x)^\lambda-1.
                                                               \tag{D.85.18}
\]
This identity holds for every \(\lambda,\nu\in\mathbb Z_5\)
by the convergent integral binomial series, not only for
ordinary nonnegative exponents.

\paragraph{The extra generator has a nontrivial tame character.}
The actual \(N_1\) multiplier and the valuation-translation
subgroup of \(I_0/N_1\) give exactly as in D.84
\[
 (J^D)^{I_0}=(T(-1)_\Gamma)^*,
 \qquad T=J/(PJ+yJ).
                                                               \tag{D.85.19}
\]
The Cartier untwist means that its compact operator is
\(\chi^{-1}\rho\). Coinvariants use the closed span of
actual operator differences.

Let \(G\) be the image of \(\mathbb Z_5g_h\) in \(T\).
This is a cyclic compact quotient, possibly zero or finite.
It is stable under the actual \(\Gamma\)-action and has
\emph{exact} character \(\chi^3\). To verify this, use
lifts with \(k_5=0\). Their compact multiplier has
augmentation \(\chi^2\); every positive augmentation term
annihilates \(g_h\) in (D.85.16). Substitution sends
\[
 g_h\longmapsto (1+g_h)^\chi-1\equiv\chi g_h
                                      \pmod{yI_L}.
\]
Thus the multiplier and substitution together give
\(\chi^3g_h\), independently of all higher actual beta
coefficients. Consequently \(G(-1)\) has character
\(\chi^2\). Its coinvariants vanish, since an actual lift
with \(\chi=2\) acts minus the identity by the unit three.

Evaluation at \(y=0\), followed by quotienting out this
generator, gives a natural exact sequence
\[
 0\longrightarrow G\longrightarrow T\longrightarrow T_0
                       \longrightarrow0,
 \qquad
 T_0=\frac{x\mathbb Z_5[[x]]}
 {f_\lambda(x)\mathbb Z_5[[x]]+
                        P(x,0)x\mathbb Z_5[[x]]}.
                                                               \tag{D.85.20}
\]
Here the product by \(P\) kills the cyclic summand of
(D.85.16), because \(P\) has zero augmentation. Equations
(D.85.16)--(D.85.18) prove exactness, including when
\(\lambda=0\), \(\nu=0\), or \(\nu\) is nonprimitive.
The map is equivariant for the actual action restricted to
\(f_{\mathbf r}=0\). It is precisely D.84's rank-one
connecting ideal with the same basis \(\mathbf u,\mathbf r\)
and local coefficient \(\lambda\).

Compact coinvariants are right exact, and the kernel just
computed has zero coinvariants. Hence the natural map gives
\[
                    T(-1)_\Gamma\simeq T_0(-1)_\Gamma.
                                                               \tag{D.85.21}
\]
This removes the extra finite-index direction integrally.
It does not discard the whole \(g_h\) term before taking
coinvariants or replace a group action by its rational rank.

\paragraph{Local depths and the actual invariant comparison.}
The three local boundary depths are determined by
\(\mathbf r\) and \(\mathbf z\). If \(d=0\), they are
all zero. If \(d>0\), exactly one boundary form, say
\(\ell\), has \(v_5(\ell(\mathbf r))=d\); its value on
\(\mathbf u\) is a unit. The identity
\[
 \ell(\mathbf z)=\lambda\ell(\mathbf u)
                       +5^h\nu\ell(\mathbf r)
\]
therefore gives, including cancellation and infinite values,
\[
 \min\{v_5(\ell(\mathbf r)),v_5(\ell(\mathbf z))\}
                         =\min(d,e).
                                                               \tag{D.85.22}
\]
The second summand cannot change this minimum; when \(d\)
is finite it has valuation at least \(h+d>d\).
All other boundary depths are zero. D.77 consequently gives
\((R^D)^{I_0}\simeq Z(\min(e,d))^*\) for \(d>0\),
and zero for \(d=0\).

D.84's full computation of the actual module \(T_0\),
including its three-jet extension and actual cyclotomic
power relation, now applies to (D.85.21). If \(e=0\),
the module \(T_0\) is zero. If \(e>0,d=0\), its
coinvariants are \(\mathbb F_5\), whereas the local
invariants are zero. If \(0<\min(e,d)<\infty\), both
invariant groups in (D.85.14) are divisible of that same
finite corank. The already proved injection between them
is an isomorphism: its dual is a surjection between free
\(\mathbb Z_5\)-modules of the same finite rank.

The case \(e=d=\infty\) requires the natural map, not
this finite-corank argument. Use the actual boundary
permutation and frame transport of D.79 to arrange
\(\mathbf r=(0,\beta)\), \(\beta\) a unit. Then
\(\lambda=0\) and \(\mathbf z=5^h\nu\mathbf r\).
The unique nonzero local invariant axis of D.77 is the
quotient \(t_2=1\), whose map is induced by the actual
surjection \(R\to\Lambda/(t_2-1)\). On this quotient,
\(g_h=0\) and \(f_{\mathbf u}\) is a unit multiple of
\(t_1-1\). The map \(J\to R\to\Lambda/(t_2-1)\)
therefore factors through exactly the evaluation map
in (D.85.20), with image the augmentation ideal.

After multiplication by \(t_1-1\) on both modules and
the Cartier untwist, this natural comparison is
\[
 I^2\hookrightarrow I,\qquad
 I=(t_1-1)\mathbb Z_5[[t_1-1]],
                                                               \tag{D.85.23}
\]
with pure substitution action. The possible extra generator
\(G\) maps to zero and has zero coinvariants by (D.85.21).
The quotient \(I/I^2=\mathbb Z_5(1)\) has nontrivial
order-four character. Exact order-four coinvariants make
(D.85.23) an isomorphism on full \(\Gamma\)-coinvariants.
Thus the natural map in (D.85.14) is an isomorphism also
at this simultaneous infinite endpoint.

We have proved the complete restriction-kernel bound
\[
 \mathcal R_s\simeq
 \begin{cases}
 \mathbb F_5,&e>0,\ d=0,\\
 0,&\text{otherwise}.
 \end{cases}                                                    \tag{D.85.24}
\]

\paragraph{The nonzero cubic class survives the finite-index global quotient.}
For the remaining case \(e>0,d=0\), the normalized cubic
cocycle \(c_3\) of D.78 has values in the constant
\(\mathbb F_5(-1)\subset B^D\). It remains nonzero in
\(H^1(I_0,B^D)\). Indeed
\[
                    B\simeq\mathbb Z_5[[y]]/(g_h).
                                                               \tag{D.85.25}
\]
Choose an actual \(\tau\in I_0\) with
\(\chi(\tau)=1,k_5(\tau)=1,c_3(\tau)=0\), as in D.84,
by correcting the cubic value with a primitive element of
\(N_1\). Its compact multiplier on (D.85.25) is
\((1+y)\mathcal B_\tau\). Since beta starts in degree
three on \(\ker\chi\), its difference from one is
\(y\) times a unit, before and after quotienting by \(g_h\).
Its fixed Cartier functionals are therefore exactly the
constant module \((B/yB)^D\).

If \(c_3\) were a coboundary, its primitive would be fixed
by \(\tau\), and hence constant. An actual element of
\(N_1\) with nonzero cubic character fixes constants and
contradicts that coboundary identity. This proves the claimed
nonvanishing. Natural coinduction and inflation in (D.85.13)
preserve it in \(H^1(G_{\mathbb Q},D)\).

Modulo five, \(\mathbf z=0\) when \(e>0\), whereas
\(\mathbf r\) is a nonboundary direction when \(d=0\).
The actual local cubic primitive of D.82 therefore applies.
Its local construction requires only these local residual
vectors, not D.82's separate global hypothesis.
At every prime away from five with nonzero abelian inertia,
D.43 and the finite-local-duality passage of D.80 give
\(H^1(G_{\mathbb Q_q},D)=0\). Only finitely many primes
have nonzero inertia in the fixed cubic coefficient module.
Enlarge the finite coefficient level to contain their local
primitives. At every other prime D.78's unramified cubic
Frobenius primitive works at its original common level.
This is D.80's finite-exception argument, and produces one
finite global coefficient module with local triviality at
all places. The global nonzero class above then shows
\(\mathcal S_s^*\ne0\). Equation (D.85.24) proves
\(\mathcal S_s=\mathbb Z/5\) in precisely this case.

\paragraph{Assembly and the actual new point.}
In the finite-index case, \(\overline{\mathbf r}\ne0\)
and \(\overline{\mathbf z}=\bar\lambda\overline{\mathbf u}\).
If \(e=0\), the local residual span has dimension two.
If \(e>0\), it is the valuation line, which is nonboundary
exactly when \(d=0\). Thus (D.85.24) and its attained
nonzero case say exactly \(\mathcal S_s=(\mathbb Z/5)^{\epsilon(W)}\).
Together with the already proved cases D.82 and D.84, this
covers every possibility under (D.85.1). The global boundary
depths are zero, and D.77's continuous splitting supplies
the remaining assertions of (D.85.4) with the exact local
sum \(n\).

For (D.85.5), \(799=17\cdot47\) and
\(800=2^5\cdot5^2\). Hence
\(L=\mathbb Z_5(1,0)+5\mathbb Z_5(0,1)\),
\(\mathbf u=(1,0)\), \(\mathbf r=(0,2)\), and \(h=1\).
The unit coefficient of \(c=\kappa(17)+\kappa(47)\)
has depth
\[
 e=v_5(799^4-1)-1=1,
 \qquad d=\infty.
                                                               \tag{D.85.26}
\]
The full rational point supplies actual section existence;
D.74 then gives the asserted entire metabelian fibre.
The supporting exact certificate
\path{code/de_sha_global_independence.py} checks the new integral
generator/action identities, the exhaustive residual case
partition, depth identities including cancellation, and
this point. The full cohomology and infinite-endpoint
claims are proved above, not inferred from those finite checks.
\end{proof}
